\section{L'interface de recherche et d'affichage}
Historiquement, CujasNum révèle que si le moteur de recherche simple s'avérait robuste et tolérant aux troncatures, l'interface souffrait quant à elle de visibilité et d'une confusion entre les différents modes de recherche. Le prototype Codex a jeté les bases de la modernisation, mais révèle des limites ergonomiques en multipliant et en répétant les onglets de recherche à différents endroits, créant une surchage cognitive. Pour surmonter ces inconvénients et répondre aux attentes formulées par des chercheurs articulent les spécificités autour de trois modalités : la recherche épurée, des filtres restructurées et des options de tri.

\subsection{La recherche simple par barre unique}
Sur CujasNum, les usagers étaient régulièrement confrontés à un bruit documentaire, ou à l'inverse à des silences inexplicables causé par la confusion entre la recherche dans les notices et la recherche plein-texte. Pour y remédier, la future page d'accueil de la bibliothèque numérique proposera une barre de recherche simple unique, positionnée au centre de l'interface. L'usager doit pouvoir sélectionner explicitement son périmètre via trois options de recherche simple : \enquote{à travers les métadonnées, la recherche s'effectue dans les notices descriptives ; à travers le plein-texte, la recherche interroge les textes océrisés et la dernière option est la combinaison des métadonnées et du plein-texte, la recherche s'effectue alors en croisant les données textuelles et descriptives de la plateforme}.

Ce modèle de recherche s'inspire des pratiques de recherche et d'autres institutions patrimoniales, comme la BIS ou l'INHA, ont développé au sein de leur propre bibliothèque numérique. Cette répartition permet d'accompagner le lecteur éliminant ainsi toute ambiguïté. Cette barre simple disposera d'un système d'autocomplétion ou de saisie intuitive, guidant l'usager dans sa recherche.

\subsection{La recherche avancée multicritères}
Dans CujasNum, le formulaire de recherche avancée était synonyme de confusion pour les agents et pour les usagers, car elle se situait \enquote{sur la même page que la recherche en texte intégrale} et utilisait un vocabulaire très professionnel, comme l'intitulé \enquote{notice bibliographique} peu compréhensible pour des étudiants ou des chercheurs non-spécialistes du livre.

La prototype Codex proposera un bouton de recherche avancée unique, clairement visibilité et situé en-dessous de la barre de recherche simple. C'est une pratique de présentation devenue courante que l'on retrouve dans différentes bibliothèques numériques, comme chez Lucienne de l'ENS ou encore Cariatide de l'INHA. Ce mode de présentation épurée mais visible réduit la charge cognitive des usagers, reproduisant des habitudes de recehrche des grands moteurs de recherche, comme Google ou Firefox. En cliquant sur le bouton de recherche avancée, l'utilisateur se trouve rediriger vers le formulaire qu'il n'a plus qu'à remplir. Ce formulaire permet ainsi de préciser des zones de recherche acceptant différents champs de saisie combiné à des opérateurs booléens.

Cette structuration permet aux chercheurs de réaliser des requêtes complexes en croisant les sources disponibles dans la bibliothèque numérique. Ce formulaire permet de répondre à l'une des attentes des chercheurs qui apprécient \enquote{de mener leur recherche à partir d'un point précis et élargissant progressivement leur périmètre de recherche au fur et à mesure}.

\subsection{Le filtrage par facettes}
L'un des chantiers identifiés au cours de l'audit de CujasNum se révèle être la gestion des filtres et des facettes. L'indexation historique reposait sur l'indexation Hauville. Cette indexatioin a généré des facettes ultra-spécifiques, souvent reliées à un seul document, comme la facette \enquote{Droit romain - Loi des Douze Tables}. Ce niveau de granularité entraînait une importante frustration pour les usagers signalant \enquote{qu'à cause de ça il n'utilisait ces facettes} ne souhaitant pas passer à côté d'un document. De surcroît, des variations de casse ou de typologies créaient davantage de facettes.

Une restructuration intégrale de ce parcours est envisagée par la bibliothèque Cujas. Le filtrage sera regroupé dans une nouvelle rubrique intitulée \enquote{Explorer}. La bibliothèque numérique proposera un ensemble de facettes générales et normalisées applicables à plusieurs types de documents. Cette rerstructuration permettra ainsi aux usagers de mieux comprendre comment sont présentés les collections. En effet, comme le démontre Leblanc \enquote{les facettes reflètent la manière dont la bibliothèque numérique a pensé ses collections et aident ainsi l'utilisateur à apprendre empiriquement l'organisation des collections}. Ce système de facettes permet ainsi la bibliothèque Cujas de guider intellectuellement les usagers. Un bon système de guidage bien formé et pensé devient un service indispensable \enquote{et permettent d'éviter le phénomène de surabondance des choix}.

Le système de filtrage ou de tri est tout aussi important que les autres modes de recherche, et nécessite d'être bien pensé. La simplification de l'indexation et des différents modes de recherche évite la confusion de l'usager, complétant ainsi les autres services proposés par la bibliothèque numérique. 